\section{Realizing a continuous law on the common bit space}
\label{sec:realization}

The normalization construction in Section~\ref{sec:continuous-laws} initially
provides an integration rule on the line. To use it as a random variable on
the common probability space, a joint realization is still needed. The
following construction supplies one. It also separates an effective sampler
in probability from a pointwise inverse-CDF algorithm. Strict positivity of
the density is unnecessary for the first construction.

\subsection{A finite monotone-coupling identity}

Let $\mu,\nu$ be finite laws with rational ordered support points and dyadic
weights. On a common sufficiently fine fair-bit table, allocate consecutive
rows to the atoms in increasing order. Denote the resulting nondecreasing
tables by $Q_\mu,Q_\nu$. Each has exactly its prescribed law. The allocations
are compatible with prefix refinement because all cumulative weights have
dyadic denominators.

\begin{lemma}[Finite monotone coupling]\label{lem:monotone-coupling}
If both laws are supported on $[-K,K]$, then
\[
 I|Q_\mu-Q_\nu|
   =\int_{-K}^{K}|F_\mu(t)-F_\nu(t)|\,dt
   \le 2K\sup_t|F_\mu(t)-F_\nu(t)|.
\]
The displayed integral is an exactly rational step-function sum.
\end{lemma}
\begin{proof}
For rational $x,y$ in the interval,
$|x-y|$ is the length of the set of thresholds where
$\mathbf1_{x\le t}$ and $\mathbf1_{y\le t}$ differ. Subdivide the threshold
interval at the finite combined list of support points and sum this identity
over the common fair-bit table. At every threshold, each of the two sets of
rows is an initial segment. Their symmetric difference therefore has mass
$|F_\mu(t)-F_\nu(t)|$. Interchange the two finite sums. The upper bound
multiplies the interval length by the largest integrand value. Sorting,
prefix sums, allocation, and this proof all use bounded rational operations.
\end{proof}

Write $T_K(x)=\max(-K,\min(K,x))$. For laws with larger support, the
same coupling satisfies
\begin{equation}\label{eq:coupling-clipped}
 d_0(Q_\mu,Q_\nu)
 \le 2K\sup_t|F_\mu(t)-F_\nu(t)|
       +\mu(|x|>K)+\nu(|x|>K).
\end{equation}
Indeed replace both tables by their clipped versions. The replacement costs
at most the two indicated exceptional masses in $d_0$; the clipped CDFs
have the same discrepancy on $(-K,K)$ and identical endpoints outside it.
Apply the lemma and $d_0\le d_1$. This estimate, rather than an unbounded
first-moment transport estimate, allows heavy-tailed laws.

\subsection{An explicit histogram sampler}

Suppose $q$ is a normalized kernel as in Section~\ref{sec:continuous-laws},
with local bounds and continuity moduli, and a PR mass-tail cutoff. Put
$\varepsilon_k=2^{-k}$. Choose a monotone rational $K_k\ge1$ such that
\[
 \int_{|x|>K_k}q(x)\,dx\le\varepsilon_k/16,
 \qquad
 \eta_k=\frac{\varepsilon_k}{64(K_k+1)}.
\]
The sequence $\eta_k$ is nonincreasing. It is essential that the following
\emph{quadrature radius} is a separate, possibly much larger number: choose
$H_k\ge K_k$ with mass tail at most $\eta_k/4$. A single radius need not
meet both budgets for a heavy-tailed law.

Partition $[-H_k,H_k]$, including all listed discontinuities, into $J_k$
rational cells of width at most $1$ and with mass at most $\eta_k/4$ each.
A supplied bound $L_{H_k}$ for the density makes any mesh of width at most
$\eta_k/[4\max(1,L_{H_k})]$ sufficient. At each cell $I_j$, quadrature
and downward dyadic rounding give a dyadic number
\[
 0\le p_j\le\int_{I_j}q,
 \qquad
 \sum_j\left(\int_{I_j}q-p_j\right)\le\eta_k/4.
\]
For example, choose a common binary denominator so that each integration
error plus rounding loss is at most $\eta_k/(4\max(1,J_k))$.
Use the cell midpoint as its atom and add an atom at zero of weight
$1-\sum_jp_j$. These are nonnegative dyadic weights summing \emph{exactly}
to one. There is no renormalization of an uncertain total mass.

Let $\mu_k$ be this finite law. Moving mass within one cell changes its CDF
by at most the mass of that cell; moving all omitted and rounded mass to
zero changes it by at most the total omitted mass. Consequently
\begin{equation}\label{eq:histogram-cdf}
 \sup_t|F_{\mu_k}(t)-F(t)|\le3\eta_k/4\le\eta_k.
\end{equation}
The comparison is proved from finite sums and then from their rational
enclosures; the supremum is just notation for the uniform inequality.
Sorted consecutive allocation turns $\mu_k$ into a cylinder table $Q_k$.

\begin{theorem}[Coherent $L^0$ realization of an effective density]
\label{thm:density-realization}
The tables $Q_k$ form a regular $L^0$ name. Their represented law is $q(x)\,dx$.
In particular the normal, exponential, and Cauchy density presentations
produce random variables on the same fair-bit probability space. Independent
copies are available for all three; conditional expectation applies to their
represented integrable observables.
\end{theorem}
\begin{proof}
For $l\ge k$, the two CDFs differ by at most $\eta_k+\eta_l$. Also
\[
 \mu_j(|x|>K_k)\le\varepsilon_k/16+2\eta_j
 \quad(j=k,l),
\]
by the two one-sided CDF comparisons. Apply
\eqref{eq:coupling-clipped} at the \emph{target} radius $K_k$:
\[
 d_0(Q_k,Q_l)
 \le 2(K_k+1)(\eta_k+\eta_l)+\varepsilon_k/8
 \le\varepsilon_k/16+\varepsilon_k/8
 <\varepsilon_k.
\]
Every table is finite and its allocation uses a declared dyadic denominator;
its prefix comparison with another table is therefore the very monotone
coupling used in the proof. This supplies coherence, not merely the existence
of a separate coupling for each pair of laws.

For a bounded Lipschitz test $h$, compare its finite expectation with the
quadrature integral. On matched cell mass the discrepancy is at most its
Lipschitz constant times the mesh width. On omitted and rounded mass it is
at most $2\|h\|_\infty$ times that mass. These errors tend to zero with
explicit PR rates. The expectation stability for $L^0$ tests then identifies
the limit with $\int hq$. This defines the law without a measure-extension
oracle. It also proves equality of the bounded-observable and density
presentations.

For a threshold $t$, the finite band probabilities satisfy
\[
 \mu_k([t-\delta,t+\delta])
 \le\int_{t-\delta}^{t+\delta}q+2\eta_k.
\]
For the left endpoint use the left limit of the finite CDF; continuity of
$F$ and rational slack removal justify the same discrepancy bound there.
A local density bound now gives the boundary certificate from
Section~\ref{sec:prob-foundation}, constructing an event name with probability
$F(t)$. Thus the threshold claim does not precede the sampler construction.
\end{proof}

If a moment-tail certificate for $|x|^p q(x)$ is supplied, this sampler also
has a corresponding integrability refinement. For $p=1,2$ it gives an
$L^p$ name after a PR reindexing. Here are the finite bounds that establish
the assertion rather than merely identify a putative moment. Since the mesh
width is at most one and each assigned cell mass is no larger than its true
mass, for $K>1$,
\[
 \mathbb E_{\mu_j}\bigl(|x|^p\mathbf1_{|x|>K}\bigr)
 \le 2^{p-1}\int_{|x|>K-1}(|x|^p+1)q(x)\,dx.
\]
The zero residual atom contributes nothing. Thus the assigned masses have
uniform effective moment tails. For $p=1$, apply the clipping/Vitali theorem.
For $p=2$, $K\ge1$ and finite tables $u,v$ give
\[
 I|u-v|^2\le12K^2d_0(u,v)+3I(R_K(u)^2)+3I(R_K(v)^2).
\]
Here $R_K(x)=(|x|-K)^+$. Insert the clipped values, use
$(a+b+c)^2\le3(a^2+b^2+c^2)$, and bound their clipped squared difference
by $4K^2\min(1,|u-v|)$. Uniform second-moment tails then give an explicit
$d_2$ Cauchy rate. More generally continuous polynomial observables have
$L^1$ names under their stated absolute-moment tails. The Cauchy identity
variable has no such first-moment refinement, as proved in the distribution
section. For exact moment agreement, first use the bounded Lipschitz test
$x\mapsto(T_Rx)^p$. Its sampler expectation equals its density integral by
the law identification already proved. The uniform moment tails on both
sides permit $R$ to grow with a PR schedule, proving
$\mathbb E Q^p=\int x^p q(x)\,dx$, and likewise for absolute moments.

\subsection{Pointwise quantiles require a different certificate}

A CDF can be continuous while having flat intervals. The preceding sampler
still works, but a unique pointwise inverse need not vary continuously in
its input data. For a safe inverse construction, suppose
\[
 F(b)-F(a)\ge c(b-a)\quad(-R\le a\le b\le R),\qquad c>0,
\]
and the target rational mass $u$ lies between $F(-R)$ and $F(R)$ with
certified margins. To locate its inverse to distance $\epsilon$, evaluate
$F(m)$ at the current midpoint to error $\delta=c\epsilon/4$. If the
enclosure is below $u$, retain the upper half; if it is above $u$, retain
the lower half. If it straddles $u$, its midpoint value has residual at most
$2\delta$ and hence spatial error at most $\epsilon/2$ by the lower-slope
inequality. Otherwise after a prescribed number of halvings the bracket
width is at most $\epsilon$. No comparison decides whether two real
numbers are exactly equal. The normal and Cauchy lower density bounds in
Section~\ref{sec:continuous-laws} supply $c$ on any certified compact range.
Existence is supplied by this construction itself: in the straddling case
intersect the current bracket with $[m-\epsilon/2,m+\epsilon/2]$; the
lower-slope estimate preserves the endpoint signs. Brackets produced at
different precisions cannot be disjoint, since strict increase would
contradict their opposite endpoint inequalities. Their midpoints are
therefore Cauchy with an explicit width bound, and continuity makes their
limit a solution. The same lower-slope estimate proves uniqueness.

\subsection{Moments, covariance, and products of integrable variables}

For $X,Y\in L^2$, the product has an $L^1$ name. In fact finite
Cauchy--Schwarz gives
\[
 \|XY-X'Y'\|_1
 \le\|X\|_2\|Y-Y'\|_2+\|Y'\|_2\|X-X'\|_2.
\]
A bound for the norm of any regular name follows from one of its finite
entries and its error. Choose input precisions with these bounds and the
displayed inequality supplies the product's PR modulus. The same argument
extends Cauchy--Schwarz. Thus
\[
 \operatorname{Var}(X)=\mathbb E(X-\mathbb EX)^2\ge0,
 \qquad
 \operatorname{Cov}(X,Y)=\mathbb E(XY)-\mathbb EX\,\mathbb EY
\]
are represented quantities. Expanding the finite squares proves all familiar
addition identities, and
$|\operatorname{Cov}(X,Y)|^2\le\operatorname{Var}(X)\operatorname{Var}(Y)$.
For every finite rational vector $(a_i)$,
$\sum_{i,j}a_i a_j\operatorname{Cov}(X_i,X_j)
=\mathbb E(\sum_i a_i(X_i-\mathbb EX_i))^2\ge0$.
This gives positive semidefiniteness of covariance matrices without importing
a classical space of random functions. An arbitrary product of two $L^1$
variables is not covered; the $L^2$ or boundedness hypotheses do real work.

There is a second useful closure rule: a represented bounded multiplier times
an $L^1$ variable is integrable, even when the multiplier converges only in
probability rather than uniformly. For finite $|a|,|b|\le M$ and $K\ge0$,
\[
 |ax-by|\le M|x-y|+K|a-b|+2M(|y|-K)^+.
\]
Use the $L^1$ error of the unbounded factor, the bounded multiplier's $L^1$
error, and the automatically supplied excess tail of $y$ to allocate these
three terms. For a bounded $L^0$ multiplier choose approximants clipped to
$[-M,M]$ and use $d_1(a,b)\le\max(1,2M)d_0(a,b)$ to obtain that $L^1$
error. This gives a PR product modulus. In particular $\mathbf1_A X$
is defined for every represented event $A$ and every represented $X\in L^1$.
This justifies the first-crossing products used in the maximal inequality.
